\begin{helper}[Disjoint two-color component normalization]
Let \((\varphi,T)\) be a minimum bad pair, write
\(T=(y_0,e_1,y_1,\ldots,e_n,y_n)\), and fix indices
\(a,b\le p<n\).  Suppose distinct colors \(\alpha,\beta\) are missing
at \(y_a,y_b\), respectively, and \(\beta\) is also missing at \(y_n\).
Let \(Q\) be the \((\alpha,\beta)\)-component rooted at \(y_n\).  If every vertex
\(y_0,\ldots,y_p\) and every exact edge-node among the protected early
edges \(e_1,\ldots,e_p\) avoids \(Q\), then the component swap
\(\varphi'=\varphi/Q\) has the following properties:
\((\varphi',T)\) is a minimum bad pair, every protected edge keeps its
exact value, and \(\alpha\) is missing at both \(y_a\) and \(y_n\).
\end{helper}
\begin{proof}
Put \(C=Q\) and let \(\varphi'\) be the exact swap from
\noderef{TwoColorComponentSwap}.  Unfold
\noderef{TashTreeMinimumBadPair} and \noderef{TashTreeBadPair} for the
original pair.  In particular, \noderef{KTriple} supplies properness of
\(\varphi\), unique uncoloredness of \(e_0=e_1\), and all the fixed
graph-and-color-count clauses.  The new coloring is proper by
\noderef{TwoColorComponentSwapProper}.  The defining cases of
\noderef{TwoColorComponentSwap} fix an uncolored value and replace a
colored value only by another colored value.  Hence the same exact edge
is uniquely uncolored, as expressed by \noderef{IsOneEdgeUncolored}, and
all remaining \(k\)-triple clauses are unchanged.

We check the Tashkinov condition edge by edge.  Fix a formal edge index
\(r\) with \(1\le r<n\), so that \noderef{TashkinovTree} supplies an old
color \(c\) of the exact edge and an old missing-\(c\) witness \(y_w\),
where \(w\le r\).  If \(r<p\), the edge-node avoids \(C\) by hypothesis,
so its value is fixed.  Also \(w\le r<p\), and hence the vertex-node
\(y_w\) avoids \(C\).  The outside-component clause of
\noderef{TwoColorComponentSwapMissingBehavior} says that its entire
missing set is fixed.  Thus the same witness proves the new Tashkinov
obligation.

It remains to treat \(p\le r\).  The cutoff inequalities give
\(a,b\le p\le r\), so both fixed witnesses are legal earlier witnesses
for the edge at index \(r\).  If \(c\notin\{\alpha,\beta\}\), the
unaffected-color clause of
\noderef{TwoColorComponentSwapMissingBehavior} preserves missingness of
\(c\) at the old witness, and the defining cases of
\noderef{TwoColorComponentSwap} preserve the edge value, so the old
witness works.  Suppose \(c=\alpha\).  If the edge-node lies in \(C\),
its new color is \(\beta\), and the fixed vertex \(y_b\), which is outside
\(C\), still misses \(\beta\).  If the edge-node lies outside \(C\), its
color remains \(\alpha\), and \(y_a\) still misses \(\alpha\).  Both
replacement indices are at most \(p\le r\), so they are legal earlier
vertices for this edge.  The case \(c=\beta\) is symmetric: use \(y_a\)
when the edge is swapped to \(\alpha\), and \(y_b\) when it stays
\(\beta\).  This proves \noderef{TashkinovTree} for \(T\) under
\(\varphi'\).

The vertex \(y_a\) avoids \(C\), so the first clause of
\noderef{TwoColorComponentSwapMissingBehavior} preserves its missing
\(\alpha\).  The root vertex-node \(y_n\) belongs to its own component;
the affected-color clause of the same result changes its missing
\(\beta\) into missing \(\alpha\).  The two vertices are distinct because
\(a\le p<n\) and the vertex map in \noderef{TreeSequence} is injective.
Consequently \(V(T)\) is not elementary under \(\varphi'\).  Together
with the new \(k\)-triple and Tashkinov condition, this is
\noderef{TashTreeBadPair}.

The sequence \(T\), hence its rank, has not changed.  The universal
minimality clause of the original \noderef{TashTreeMinimumBadPair}
quantifies over every coloring, so it combines with the new badness to
make \((\varphi',T)\) minimum bad.  Finally, every protected edge-node
is outside \(C\), and its exact value is therefore fixed directly by
\noderef{TwoColorComponentSwap}.  This proves all asserted conclusions.
\end{proof}
